\section{Preliminaries}\label{sec:prelims}

As mentioned, we assume the reader is familiar with the basics of parameterized
complexity, such as the class FPT~\cite{cygan2015parameterized}, and approximation algorithms
such as a PTAS~\cite{williamson2011design}. A
\emph{parameterized approximation scheme (PAS)} is an algorithm that computes a
$(1+\eps)$-approximation for a problem in  $f(k,\eps)n^{g(\eps)}$ time for some
functions $f$ and $g$, while an an \emph{efficient parameterized approximation
scheme (EPAS)} is a $(1+\eps)$-approximation algorithm running in time
$f(k,\eps)n^{O(1)}$ (that is, the running time is FPT in $k+\frac{1}{\eps}$).
The distinction between a PAS and an EPAS is similar to the one between a PTAS
and an EPTAS.

By $w:E\to\mathbb{R}^+$ we denote an edge-weight function, so that the cost of a solution $F$
to the \SF problem is $\cost(F)=\sum_{e\in E(F)} w(e)$. 
We will use $F^\star$ to denote an optimal solution, and for
$\alpha\ge 1$ we will say that a solution $F$ is $\alpha$-approximate if
$\cost(F)\le \alpha\cost(F^\star)$. For $u,v\in V$ we use $\dist(u,v)$ to
denote the shortest-path distance from $u$ to $v$ in $G$ according to the weight function~$w$.

\begin{definition} Given a graph $G=(V,E)$, a \emph{tree decomposition} is a
pair $(T,\{B_i\}_{i\in V(T)})$, where $T$ is a tree and each node $i\in V(T)$
of the tree is associated with a \emph{bag} $B_i\subseteq V$, with the
following properties: 
\begin{enumerate}
\item $\bigcup_{i\in V(T)} B_i=V$,
i.e., all vertices of $G$ are covered by the bags, 
\item for every edge $uv\in
E$ of $G$ there exists a node $i\in V(T)$ of the tree for which $u,v\in B_i$,
and 
\item for every vertex $v\in V$ of $G$ the nodes $\{i\in V(T)\mid v\in
B_i\}$ of the tree for which the bags contain $v$ induce a (connected) subtree
of $T$.  
\end{enumerate} 
The \emph{width} of the tree decomposition is
$\max_{i\in V(T)}\{|B_i|-1\}$ and the \emph{treewidth} of $G$ is the minimum
width over all its tree decompositions.  

A rooted tree decomposition is \emph{nice} if for every $i\in V(T)$ we have one of the
following: 

\begin{enumerate} 

\item $i$ has no children\footnote{here we do not demand the leaf nodes to be
empty, as is often assumed for this definition.} ($i$ is a \emph{leaf node}),

\item $i$ has exactly two children $i_1$ and $i_2$ such that
$B_i=B_{i_1}=B_{i_2}$ ($i$ is a \emph{join node}), 

\item $i$ has a single child $i'$ where $B_i=B_{i'}\cup \{v\}$ for some $v\in
V$ ($i$ is an \emph{introduce node}), or 

\item $i$ has a single child $i'$ where $B_i=B_{i'}\setminus \{v\}$ for some
$v\in V$ ($i$ is a \emph{forget node}).  

\end{enumerate}

\end{definition}

Given a rooted tree decomposition $T$ of a graph $G$, for a node $u$ of $T$ let
$B$ be the bag associated with it.  Then $V_B$ is the set of vertices of all
bags in the subtree rooted at $u$.  The set $A_B\subseteq R$ denotes the
\emph{active terminals} of the bag $B$: for any demand pair~$\{s,t\}\in D$, if
$s\in V_B$ and $t\notin V_B$ then $s\in A_B$.  For any \SF solution $F$, if a
connected component $C$ of $F$ contains an active terminal, then we say that
$C$ is an \emph{active component} for $B$. For a fixed solution $F$, we denote
the set of all active components for $B$ by~$\mc{C}_B$.  Note that every active
component must intersect the bag $B$. 

If for every bag $B$ a set of partitions $\Pi_B$ of $A_B$ is given, a
\SF solution $F$ is \emph{conforming} to all $\Pi_B$, if for each bag~$B$ there exists a partition
$\pi\in\Pi_B$ such that any two active terminals in $A_B$ are in the same set
$S\in\pi$ if and only if they are part of the same active component $C$ of~$F$,
i.e., $S\subseteq V(C)$ and $S'\cap V(C)=\emptyset$ for any $S'\in\pi$ with
$S'\neq S$ (note that this implies $|\pi|\leq |B|$). One technicality of \cref{thm:tw-dp} is that
the algorithm needs a nice tree decomposition as input, for which the terminals
only appear in bags that are leaf nodes of the decomposition. Given any tree
decomposition, these conditions are not hard to meet
(cf.~\cite[Lemma~6]{bateni2011approximation}).  However, for our algorithms we
are going to rely on tree decompositions with certain additional properties.
Hence we will need to revisit the conditions needed for the algorithm of
\autoref{thm:tw-dp} when using it for our purposes.

We will also consider the following parameters: 
The \emph{treedepth} of a graph $G$ can be defined recursively as follows: (i) the treedepth of $K_1$ is $1$ (ii) the treedepth of a disconnected graph is the maximum of the treedepth of any of its components (iii) the treedepth of a connected graph $G$ is $1+\min_{v\in V(G)} \textrm{td}(G-v)$. 
A \emph{feedback vertex set} is a set of vertices whose removal leaves a forest. A \emph{vertex cover} is a set of vertices such that its removal leaves an edge-less graph. A \emph{feedback edge set} is a set of edges whose removal leaves a forest. In a connected graph with $n$ vertices and $m$ edges the minimum feedback edge set always has size $m-n+1$.

As part of our approximation algorithm we will use the notion of $\delta$-nets, defined as follows. 
A well-known fact is that a $\delta$-net exists for any metric and any $\delta\geq 0$, 
and it can be constructed greedily in polynomial time. 

\begin{definition}\label{dfn:net}
Given a metric $(X,\dist)$, a \emph{$\delta$-net} is a subset $N\subseteq X$ of points, such that
\begin{enumerate}
    \item any two net points $u,v\in N$ are far from one another, i.e., $\dist(u,v)>\delta$, and
    \item for any node $u\in X$ there is some net point $v\in N$ close by, i.e., $\dist(u,v)\leq\delta$.
\end{enumerate}
\end{definition}

